% !Mode:: "TeX:UTF-8"
\chapter{绪论}
\label{cha:chap01}

\section{研究背景与意义}

细胞是生命活动的基本结构与功能单元。细胞群体在基因表达、代谢状态、形态结构、迁移能力和力学表型等方面具有明显差异，群体平均信号难以完整反映单个细胞的状态分布。单细胞测序、质谱流式细胞术、基于液滴条形码的单细胞转录组测序和高通量成像等技术推动了细胞状态分析由群体统计向单细胞分辨发展\cite{altschuler2010heterogeneity,trapnell2015singlecell,shapiro2013singlecell,bendall2011singlecellmass,macosko2015dropseq,klein2015dropletbarcoding,heath2016singlecelltools,regev2017humancellatlas}。单细胞力学表型与分子信息和形态信息相互补充，为描述细胞在外部力学环境中的状态及响应提供了物理维度。

细胞力学行为由细胞膜、膜皮质层、细胞骨架、胞质和细胞核等结构共同决定。细胞骨架承担力的传递与重构，细胞膜和膜皮质层维持细胞边界与整体形态，胞质流动和细胞核变形参与时间相关的变形与恢复过程\cite{kasza2007cellmaterial,fletcher2010cellmechanics,wang2009mechanotransduction,moeendarbary2013poroelastic}。基质刚度能够调节细胞黏附、迁移和分化行为\cite{discher2005tissue,engler2006matrix}。肿瘤细胞、血细胞和免疫细胞的弹性、黏弹性及变形能力也会随细胞类型和状态发生变化，力学表型因而成为细胞状态表征中的一类重要指标\cite{suresh2007cancer,dicarlo2012mechanical,toepfner2018morphorheological}。

细胞弹性模量、黏性系数、黏弹性特征时间、膜张力和整体变形能力通常不能由静态图像直接获得，而需要在外部载荷和边界条件已知的情况下，根据力与位移关系、形态变化、时间响应或恢复过程进行估计。原子力显微镜、微吸管吸吮、光学拉伸器和磁镊等方法能够获得局部或整体细胞力学响应，但通常依赖逐细胞定位、局部加载或较复杂的实验操作。高通量细胞力学测量需要在连续处理大量细胞的同时保持加载条件、观测量和参数解释之间的一致性\cite{darling2015high,dicarlo2012mechanical,urbanska2020comparison}。

细胞形态与细胞力学参数之间存在联系，但二者并不等同。静态形态会受到细胞类型、细胞周期、渗透状态和成像姿态影响，外形相近的细胞可以具有不同的材料性质。外部载荷作用后的动态形态记录变形建立、峰值变形和恢复过程，其中变形幅值与材料刚度相关，响应滞后和恢复速度包含黏弹性信息，质心运动还受到流体载荷和几何阻力影响\cite{mietke2015rtdc,mokbel2017rtdc,dudani2016viscoelastic,gong2024prediction}。因此，动态形态为力学参数估计提供了比单帧静态几何量更丰富的过程信息。

微流控通道能够通过压力驱动流动、局部收缩、剪切流和伸展流对细胞施加可控载荷，并在细胞运动过程中连续记录轮廓和位置。收缩通道中的细胞依次经历进入约束区、持续变形、通过最小截面和离开约束后的恢复过程。细胞轮廓、质心位置、通过时间和运动速度共同记录受限变形的时空演化。实时变形细胞术和定量变形细胞术显著提高了细胞形态和力学表型的测量通量\cite{otto2015rtdc,nyberg2017qdc,rosendahl2018rtdc}。微流控形态响应同时受到流体载荷、通道几何、细胞尺寸和材料参数影响，因而需要结合具体力学模型解释形态与材料参数之间的关系\cite{mietke2015rtdc,mokbel2017rtdc}。

实际微流控成像会受到帧率、视场范围、曝光条件和图像分割误差影响。完整形态序列可能被压缩为少量观测时刻，局部轮廓也可能受到边界噪声干扰。单帧形态描述特定位置处的变形状态，完整序列还包含变形建立、峰值变形和形态恢复的过程信息。多时刻轮廓、细胞位置和时间信息可用于弹性模量预测\cite{gong2024prediction}。在稀疏观测条件下，形态恢复不仅需要保持轮廓几何，还需要保留后续力学参数估计所依赖的动态信息。

含细胞复合液滴下降接触构成另一类瞬态流体加载过程。液滴接触固体壁面后经历快速铺展、最大铺展、回缩、反弹或沉积，外界面形态和接触线运动由惯性、黏性、表面张力和壁面润湿性共同控制\cite{rein1993impact,yarin2006impact,josserand2016drop,snoeijer2013moving}。液滴内部包含细胞时，壁面接触引起的压力和速度场变化通过包裹液体传递至细胞边界，细胞变形和位移又会改变局部流动阻力与能量分配。含细胞复合液滴由此能够同时产生液滴铺展、动态接触角、细胞轮廓和质心运动等多类响应。Tasoglu 等采用复合液滴模型分析含细胞液滴沉积及内部变形\cite{tasoglu2010compound}，复合液滴撞击研究也表明内部相会改变外部液滴的铺展、回缩和界面运动\cite{blanken2021compoundperspective}。

微通道细胞受限变形与含细胞复合液滴下降接触均由流体载荷驱动细胞产生动态形态响应。前者具有稳定的几何边界和压力条件，适合研究连续受限加载下的轮廓演化与参数反演。后者同时包含自由界面、表面张力、壁面润湿和内部尺寸等因素，适合研究瞬态多相加载下的内外形态耦合和多参数响应。两类场景在物理边界和加载历程上存在明显差异，但均可按照正问题生成动态形态，再由观测形态估计材料参数的基本关系进行研究。

\section{细胞力学特性的测量与形态响应反演}

细胞力学测量由外部加载、响应观测和参数计算三个环节组成。传统方法通过探针、吸管、光场或磁场施加局部或整体载荷，微流控方法通过连续流动和微尺度几何诱导细胞变形。两类方法在加载范围、测量通量、观测时间尺度和参数解释方式方面存在差异。形态响应反演则以细胞轮廓、位移和时间演化为观测量，在给定力学模型或数据驱动映射下估计材料参数。

\subsection{传统单细胞力学测量方法}

原子力显微镜（atomic force microscopy，AFM）通过悬臂梁探针与细胞表面接触，记录加载力与压入深度之间的关系。Binnig 等提出 AFM，为微纳尺度表面形貌和相互作用力测量建立了仪器基础\cite{binnig1986afm}。细胞力学实验通常将压入曲线与赫兹模型或其修正形式结合，计算细胞表观弹性模量。探针形状、压入深度、加载速率、细胞黏附状态和基底效应均会改变模量估计结果\cite{gavara2016afm}。AFM 适合测量细胞局部刚度及其空间异质性，也可识别细胞骨架药物处理或病理状态变化对应的力学差异\cite{rotsch2000cytoskeletal,cross2007nanomechanical}。

微吸管吸吮通过负压将细胞局部吸入微吸管，并记录吸入长度随压力和时间的变化。薄壳、液滴和黏弹性模型可用于计算膜张力、皮质张力、表观黏度和整体变形参数。Hochmuth 系统总结了活细胞微吸管吸吮中的力学模型和实验处理方法\cite{hochmuth2000micropipette}。Evans 和 Yeung 根据血液粒细胞的吸入过程计算表观黏度和皮质张力，体现了该方法对整体黏弹响应的表征能力\cite{evans1989granulocytes}。微吸管吸吮适合悬浮细胞以及细胞膜和膜皮质层的力学测量，吸管尺寸、入口形状、接触条件和初始细胞形态会影响吸入过程及参数结果。

光镊和光学拉伸器利用光场对细胞或微粒施加载荷。单束梯度力光阱为微尺度非接触操控提供了基础\cite{ashkin1986singlebeam}。光学拉伸器采用两束相向激光对悬浮细胞施加整体拉伸，通过长轴和短轴变化及恢复过程描述细胞变形能力\cite{guck2001optical}。光学变形性可用于区分不同细胞状态\cite{guck2005opticaldeformability}。光学方法避免了机械探针直接接触，但有效载荷分布与细胞折射率、激光功率和热效应有关。光镊还可通过捕获微球或细胞局部结构测量细胞膜、细胞骨架和黏附相关作用力\cite{svoboda1994optical}。

磁镊和粒子跟踪微流变主要用于局部黏弹性和细胞内部微环境测量。磁镊通过细胞表面结合磁珠施加力或力矩，局部位移用于计算细胞骨架和细胞膜附近的黏弹响应\cite{bausch1999magnetic}。活细胞微流变实验显示，细胞在较宽时间尺度内可呈现幂律流变特征\cite{fabry2001scaling}。粒子跟踪微流变根据示踪颗粒的热运动或主动响应估计胞质和细胞骨架网络的局部流变性质\cite{wirtz2009microrheology}。传统单细胞力学方法具有较明确的加载边界和较强的参数解释能力，但逐细胞定位和重复加载限制了大样本测量效率。

\subsection{基于微流控的高通量细胞力学测量方法}

微流控系统利用微米尺度通道控制压力、速度梯度、流线结构和几何约束。低雷诺数流动使黏性作用和边界条件对细胞运动具有直接影响，通道几何可重复地限定细胞受力位置和变形路径\cite{squires2005microfluidics,stone2004engineering}。细胞在连续流动中依次通过加载区和观测区，高速成像同步记录面积、圆度偏差、长宽比、轮廓和质心位置。Lincoln 等提出变形性流式细胞术，将细胞形态变化与流式检测结合\cite{lincoln2004deformability}。Tse 等利用单细胞力学表型开展病理样本分析，说明形态与流变信息可用于细胞状态识别\cite{tse2013malignant}。

收缩通道方法通过窄缝、微孔或渐缩结构强制细胞发生整体变形。细胞进入受限区域后，轴向长度、径向宽度、通过速度和停留时间随材料性质和几何限制发生变化。Shelby 等建立感染疟原虫红细胞通过毛细通道的微流控模型，展示了细胞变形能力与微循环阻塞之间的关系\cite{shelby2003capillary}。基于变形和表面摩擦的微流控测量可区分不同癌细胞力学表型\cite{byun2013deformability}，微流控棘轮结构还可依据细胞变形能力进行分选\cite{park2016ratchets}。收缩通道能够产生较大的整体变形，同时需要考虑壁面接触、摩擦、黏附和堵塞对响应的影响。

伸展流和剪切流方法由流场本身对细胞施加非接触载荷。Gossett 等利用流动聚焦区域的伸展流对单细胞进行水动力拉伸，实现大样本力学表型测量\cite{gossett2012hydrodynamic}。Dudani 等在微流控伸展流中测量细胞随时间变化的变形，并利用加载和恢复过程表征黏弹性\cite{dudani2016viscoelastic}。这类方法减少了细胞与固壁的直接接触，其响应仍取决于速度梯度、细胞尺寸、初始位置和流体黏度。

实时变形细胞术（real-time deformability cytometry，RT-DC）利用微通道内的水动力应力使细胞连续变形，并在流动过程中实时提取面积、圆度偏差和变形量。Otto 等报道的 RT-DC 可在较高通量下获得单细胞力学表型\cite{otto2015rtdc}。Rosendahl 等将一维荧光成像与实时变形测量结合，使形态、力学和分子标记信息同步获得\cite{rosendahl2018rtdc}。Mietke 等结合低雷诺数流动和线弹性理论，建立由细胞尺寸和变形量计算细胞刚度的关系，并采用力学性质已知的模型球及生物细胞进行验证\cite{mietke2015rtdc}。图~\ref{fig:rtdc_mietke2015}给出了 RT-DC 中的通道结构、观测位置和随细胞运动的坐标描述。

\begin{figure}[htbp]
  \centering
  \includegraphics[width=0.82\textwidth]{figures/chap01/fig_rtdc_mietke2015.png}
  \caption[实时变形细胞术中的通道结构与受限流动示意图]{实时变形细胞术中的通道结构与受限流动示意图。图（a）给出了微通道几何、细胞变形区域和图像采集位置。图（b）给出了随细胞运动的坐标系及细胞周围的流动结构。引自文献\cite{mietke2015rtdc}。}
  \label{fig:rtdc_mietke2015}
\end{figure}

定量变形细胞术侧重测量标定和不同实验之间的可比性。Nyberg 等通过标定化流体加载获得细胞力学性质\cite{nyberg2017qdc}。Urbanska 等比较了收缩通道、剪切流和伸展流等高通量细胞变形方法，指出各方法在加载机制、通量、壁面接触、可测参数和模型依赖性方面具有不同适用范围\cite{urbanska2020comparison}。微流控方法提高了形态数据获取效率，但将面积、圆度偏差、长宽比和通过时间等观测量转换为弹性模量或黏弹性参数，仍需给出流体载荷、通道几何、细胞模型和参数标定方式。

\subsection{微通道细胞变形与力学参数反演}

细胞在变截面微通道中运动时，依次经历进入收缩段、持续变形、通过狭窄区域和离开约束后的恢复过程。入口附近的轮廓主要反映初始形态和流动扰动，收缩段内的轮廓记录压力、黏性剪切和几何限制共同作用下的变形，出口后的轮廓和速度反映材料恢复与细胞运动。完整过程中的轮廓、位置和时间信息构成由形态响应估计材料参数的数据基础\cite{mietke2015rtdc,mokbel2017rtdc,gong2024prediction}。

微通道形态的力学解释依赖流固耦合模型。细胞受到流体压力、黏性剪切、壁面约束和材料恢复力共同作用，其变形不仅由弹性模量决定，还与细胞尺寸、泊松比、黏性参数及通道几何有关。Mokbel 等采用数值方法研究 RT-DC 中细胞力学性质的提取，比较不同材料模型和皮质层参数对细胞形态的影响\cite{mokbel2017rtdc}。图~\ref{fig:rtdc_mokbel2017}显示，小变形条件下不同模型得到的轮廓较为接近，变形幅度增大后，本构关系和三维形态差异对二维观测轮廓的影响逐渐增强。

\begin{figure}[htbp]
  \centering
  \includegraphics[width=0.90\textwidth]{figures/chap01/fig_rtdc_mokbel2017.png}
  \caption[不同力学模型和参数条件下的 RT-DC 细胞形态预测结果]{不同力学模型和参数条件下的 RT-DC 细胞形态预测结果。上排比较不同材料模型和参数条件下的细胞轮廓，下排给出三维细胞后端形态及其侧向观测差异。引自文献\cite{mokbel2017rtdc}。}
  \label{fig:rtdc_mokbel2017}
\end{figure}

复杂通道中的细胞变形可通过数值模拟建立材料参数与形态响应之间的数据库。可变截面通道使细胞在有限距离内经历连续变化的流体载荷和几何约束，参数变化对应的轮廓差异能够在多个位置被记录。Xie 和 Hu 的可变截面微通道模型分析了胶囊通过局部收缩区域时的边界变形和门控行为，说明膜力学参数与通道形状共同决定通过过程\cite{xie2021capsule}。解析反演适用于模型形式和边界条件较明确的情况，材料大变形、黏弹性、壁面接触和多阶段运动则需要数值标定、迭代优化或数据驱动映射。Gong 等将多个时刻的细胞轮廓、时间和空间位置信息输入卷积神经网络，用于预测弹性模量和材料本构类别\cite{gong2024prediction}，说明多时刻形态能够提供单帧几何指标未覆盖的过程信息。

完整轮廓序列的获取受到成像帧率、曝光时间、视场范围和分割误差限制。观测帧数减少时，峰值变形或恢复阶段可能缺失。边界噪声会改变局部曲率、周长和高频轮廓特征。稀疏形态恢复需要同时评价局部边界、整体几何和动态过程。点坐标表示能够保留局部轮廓细节，但维度随采样点数增加，并对起始点、点序和局部扰动较为敏感。傅里叶轮廓表示将闭合边界转换为有限阶系数，以较低维度描述整体尺度、主变形方向和局部变化\cite{kuhl1982elliptic}。形态、质心位置和速度的联合建模还可将轮廓变形与细胞运动联系起来\cite{caruana1997multitask}。

\section{含细胞复合液滴下降接触过程及其数值模拟研究}

含细胞复合液滴由外部包裹液体、内部细胞和周围介质组成。液滴下降并接触固壁后，外部自由界面发生铺展和回缩，内部流体由轴向运动转为径向运动，细胞在压力梯度、黏性剪切和几何约束作用下发生变形与迁移。外部液滴和内部细胞具有不同的响应尺度，二者通过内部流场和界面运动相互影响\cite{tasoglu2010compound,blanken2021compoundperspective,freund2014bloodcells}。该过程的观测量既包括铺展因子、接触线位置和动态接触角，也包括细胞轮廓、质心位置、速度和恢复过程。

\subsection{液滴下降接触与动态润湿}

单一液滴下降接触固体壁面后，可经历接触初期、快速铺展、最大铺展、回缩、反弹或沉积等阶段。Rein 对液滴接触固体和液体表面的典型形态进行了分类\cite{rein1993impact}。Yarin 从铺展、回缩、飞溅和反弹等方面总结了液滴撞击动力学\cite{yarin2006impact}。Josserand 和 Thoroddsen 归纳了惯性、黏性、表面张力、周围气体和壁面性质对液滴撞击结果的共同影响\cite{josserand2016drop}。

液滴接触壁面的早期阶段主要受惯性作用控制。轴向动量在近壁区域转化为径向动量，液滴底部形成向外扩展的薄液层，接触线快速移动。随着径向速度降低，新生成的气液界面、近壁黏性耗散和毛细恢复逐渐影响铺展过程。Clanet 等研究低黏度液滴的最大变形标度，指出惯性和毛细作用控制最大铺展尺度\cite{clanet2004maximal}。Mao 等基于能量守恒分析液滴铺展和反弹，给出了黏度、接触角和冲击条件对最大铺展的影响\cite{mao1997spread}。Laan 等将边界层耗散和毛细限制纳入最大直径关系\cite{laan2014maximumdiameter}，Wildeman 等则从能量预算角度分析了撞击铺展中的几何损失和黏性耗散\cite{wildeman2016spreading}。

Pasandideh-Fard 等通过实验和数值模拟分析液滴撞击固壁过程中的毛细效应，并将动态接触角用于接触线边界处理\cite{pasandidehfard1996capillary}。图~\ref{fig:drop_impact_pasandideh1996}给出了水滴撞击不锈钢表面后的典型形态序列。完整时间序列同时记录铺展速率、最大半径、前缘液体积聚和回缩路径，其信息量高于单一最大铺展值。

\begin{figure}[htbp]
  \centering
  \includegraphics[width=0.86\textwidth]{figures/chap01/fig_drop_impact_pasandideh1996.png}
  \caption[水滴撞击不锈钢表面过程中的实验图像与数值结果]{水滴撞击不锈钢表面过程中的实验图像与数值结果。液滴在惯性作用下沿壁面铺展，达到最大铺展状态后在表面张力和黏性作用下回缩。引自文献\cite{pasandidehfard1996capillary}。}
  \label{fig:drop_impact_pasandideh1996}
\end{figure}

壁面润湿性通过接触角和接触线运动改变铺展与回缩。静态接触角描述平衡状态下固液、固气和液气界面能之间的关系，动态接触角随接触线速度和局部流动状态变化。移动接触线附近存在由微观尺度到宏观尺度的跨尺度问题，滑移、界面曲率和黏性耗散均会影响动态润湿过程\cite{deGennes1985wetting,bonn2009wetting,snoeijer2013moving}。铺展阶段的前进接触角和接触线速度共同影响径向扩展，回缩阶段的后退接触角、接触角滞后和接触线钉扎则影响液滴的回缩、反弹和沉积状态\cite{quere2008wettingroughness,bird2013contact}。

\subsection{复合液滴内部载荷传递与流固耦合}

复合液滴包含外部自由界面和内部组分边界。双乳液和核壳液滴研究表明，内部相的位置、尺寸、黏度和界面性质会改变液滴生成、输运和整体稳定性\cite{teh2008droplet,seemann2012droplet,utada2005doubleemulsions,chu2007multipleemulsions}。含细胞复合液滴中的内部细胞具有有限尺寸和可变形性，细胞边界能够储存弹性能并产生黏弹性耗散，其力学响应不能由液液界面模型完整描述。

Tasoglu 等采用前沿追踪方法研究复合黏性液滴接触平面的过程，将内部液滴用于近似被包裹细胞，并分析雷诺数、尺寸比、界面张力比、黏度比和接触角对内部变形的影响\cite{tasoglu2010compound}。Blanken 等总结了核壳和其他复合液滴撞击研究，指出内部相会改变铺展、飞溅、回缩和自润滑行为，实验测试和数值模拟均需同时处理多个界面及其时间演化\cite{blanken2021compoundperspective}。这些研究构成了含细胞复合液滴内外响应耦合分析的基础。

液滴接触壁面后，外部流动和内部细胞之间的载荷传递具有明显阶段性。接触前，细胞主要随液滴整体下降。初始接触后，液滴底部压力迅速升高，轴向速度降低，内部流体转为径向运动。快速铺展后，内部回流和界面恢复继续改变细胞表面的压力和剪切分布。细胞弹性影响瞬时变形幅度和恢复趋势，黏弹性影响响应时序和能量耗散，细胞尺寸则改变与液滴界面及壁面之间的液体间隙\cite{freund2014bloodcells,barthesbiesel2016motion}。

流固耦合模型同时求解流体运动和细胞变形。流体在细胞边界处传递法向压力和切向黏性作用，细胞位移改变流体域几何和局部阻力。细胞流动数值研究表明，材料参数、内外黏度比和流场类型共同决定细胞的变形、迁移和应力分布\cite{freund2014bloodcells,barthesbiesel2016motion}。浸没边界方法通过欧拉流体网格与拉格朗日结构边界之间的力和速度交换处理流固耦合\cite{peskin2002immersed}。任意拉格朗日欧拉方法采用随结构运动的计算网格保持流固界面清晰，适用于界面可跟踪且拓扑不发生破裂的变形过程\cite{donea2004ale}。

\subsection{多源形态响应与参数可辨识性}

含细胞复合液滴下降接触可产生外部液滴和内部细胞两组动态观测。外部观测包括铺展因子、液滴高度、接触线位置、动态接触角和回缩过程，内部观测包括细胞长宽比、圆度偏差、质心位移、速度、最大变形时刻和恢复过程。不同响应量对参数的敏感阶段和敏感方向并不相同\cite{tasoglu2010compound,blanken2021compoundperspective,dudani2016viscoelastic}。

材料参数首先改变细胞自身的变形和恢复。弹性模量较低时，细胞在相同流体载荷下产生较大形变。弹性模量较高时，边界偏离初始形态的程度较小。黏弹性特征时间改变加载与响应之间的时间关系，使峰值变形时刻、恢复速度和滞后程度发生变化。流动、界面和几何参数则直接改变输入载荷或载荷传递路径。下降速度影响惯性作用和铺展强度，表面张力影响外部界面的恢复能力和毛细时间尺度，壁面接触角影响接触线运动，细胞尺寸改变内部对象受力面积、液体间隙和几何阻塞程度\cite{tasoglu2010compound,blanken2021compoundperspective,josserand2016drop}。

参数敏感性与参数可辨识性具有不同含义。某一响应随参数变化而明显改变，表示该响应对参数敏感。不同参数产生相近响应时，仅凭该响应仍难以区分参数来源。可辨识性取决于参数与响应之间的映射是否具有足够差异，也受到观测噪声、时间采样和参数范围影响\cite{raue2009identifiability}。最大变形可同时受到弹性模量和下降速度控制，最大铺展可同时受到速度、表面张力和接触角控制。峰值时刻、恢复过程和内外响应之间的相对时序能够补充幅值指标未包含的信息。

\subsection{含细胞复合液滴的数值模拟方法}

液滴下降接触的数值模拟需要描述自由界面大变形和移动接触线。体积分数法通过网格单元中的体积分数表示不同流体占比，具有较好的质量守恒性质，界面曲率和表面张力计算依赖界面重构\cite{hirt1981vof}。水平集方法采用连续符号距离函数隐式表示界面，便于计算法向和曲率，并需要进行质量守恒修正\cite{osher1988fronts,sussman1994levelset}。前沿追踪方法显式维护界面网格，可保持较清晰的界面位置，大变形和拓扑变化时需要重构界面\cite{unverdi1992fronttracking}。

相场方法以有限厚度的弥散界面替代尖锐界面，并通过自由能和相场变量描述界面演化。弥散界面方法可将表面张力写入连续体方程，并与两相不可压缩流动方程耦合\cite{anderson1998diffuseinterface,jacqmin1999phasefield,badalassi2003phasefield}。壁面自由能或相应边界条件用于给定平衡接触角，动态接触线行为还受到界面厚度、迁移率、网格和时间步影响。细胞边界可采用膜、壳或连续体实体描述，模型选择决定材料参数的物理含义和可输出变量\cite{freund2014bloodcells,barthesbiesel2016motion}。

数值结果的可靠性包括模型验证和离散误差控制。单一液滴的铺展过程、最大直径和接触角演化可用于检验自由界面与润湿模型，可变形结构在流体载荷下的变形可用于检验材料本构和流固耦合。网格和时间步分析需要覆盖外部铺展、内部变形、质心运动和应力等主要输出\cite{roache1998verification}。当数值模型生成的形态序列用于参数反演时，离散误差、轮廓提取方式和时间采样规则也会进入训练数据。

\section{人工智能在形态重构与力学参数反演中的应用}

形态响应数据可表示为图像、轮廓点、几何指标和时间序列。形态重构以稀疏或含噪观测为输入，输出完整轮廓或连续形态过程。参数反演以一个或多个时刻的形态响应为输入，输出弹性模量、黏弹性参数或其他待识别参数。两类任务共享特征提取和时序建模方法，但评价目标不同。重构任务关注几何保真和动态连续性，参数反演关注材料参数误差和不同观测条件下的稳定性。

\subsection{数据驱动的形态表征与参数预测}

早期数据驱动材料建模主要使用多层前馈神经网络拟合应力与应变之间的关系或材料状态变量之间的映射。Ghaboussi 等将神经网络用于材料行为表示，展示了由数据学习非线性材料响应的基本形式\cite{ghaboussi1991materialnn}。数据驱动计算力学将材料数据与平衡条件、相容条件和边界条件联系起来，在状态空间中寻找满足力学约束的材料响应\cite{kirchdoerfer2016datadriven}。细胞形态参数反演还需要处理高维边界、图像和时间序列。

多层感知机适合处理固定长度向量，可将若干观测时刻的轮廓坐标或几何指标拼接后进行回归。卷积神经网络利用局部卷积核提取相邻点或图像区域的空间特征，适用于规则采样轮廓和二维图像。残差网络等结构为高维特征提取提供了通用方法\cite{lecun2015deeplearning,krizhevsky2012imagenet,he2016resnet}。细胞轮廓属于闭合曲线，卷积边界处理、起始点和点序规则会影响首尾邻接关系。

变分自编码器将高维输入编码为概率潜变量，并由潜变量重构原始数据\cite{kingma2014vae}。该方法能够学习轮廓的低维分布并生成缺失形态。Transformer 采用注意力机制建模序列中不同位置之间的依赖关系\cite{vaswani2017attention}。在稀疏形态重构中，观测时刻可表示为序列元素，时间和空间位置通过位置编码进入模型。注意力机制能够建立相距较远时刻之间的联系，轮廓闭合和几何量一致性则需要通过数据处理、网络结构或损失函数进一步保证。

轮廓输入可分为点坐标、参数曲线和几何指标。点坐标保留局部边界细节，维度较高且依赖轮廓点对应关系。面积、圆度偏差和长宽比等几何指标维度低且解释直接，但会丢失局部形态和边界方向信息。傅里叶轮廓表示利用有限阶谐波描述闭合边界，能够压缩轮廓并控制重构精度\cite{kuhl1982elliptic}。傅里叶阶数过低会产生欠拟合，阶数过高则会重新引入边界噪声。图像输入保留外部液滴、内部细胞和壁面接触区域的空间关系，同时也包含像素分辨率、裁剪范围和预处理规则等观测因素。

时间信息可以通过多帧拼接、时序卷积或注意力机制处理。单帧输入适合评价特定阶段的参数敏感性，多帧输入能够同时利用加载和恢复过程。依据物理阶段选取接触前、最大压缩、最大铺展和恢复等代表性时刻，可减少相邻帧的重复信息。多任务学习可同时输出轮廓、质心速度和材料参数\cite{caruana1997multitask}。共享编码器提取动态形态特征，不同输出分支分别处理几何、运动和参数任务，几何派生量和质心运动关系可作为多任务学习中的一致性约束。

\subsection{物理先验、模型约束与本构参数反演}

数据驱动模型以样本误差为主要训练依据，控制方程、本构形式、材料对称性、能量关系、几何边界和已知变量关系可以通过损失函数、网络结构或输入变量进入学习过程。物理信息神经网络在损失函数中加入偏微分方程残差、初始条件和边界条件，可用于正问题求解和参数反演\cite{raissi2019pinn}。广义物理信息机器学习还包括算子学习、多保真数据和物理结构编码\cite{karniadakis2021piml}。

材料本构学习中的约束集中于客观性、材料对称性、能量关系和耗散非负。热力学人工神经网络将自由能和耗散关系写入网络结构或训练目标，使预测结果满足热力学条件\cite{masi2021tann}。Flaschel 等利用位移场和全局力数据，在平衡约束和稀疏正则化下识别可解释的超弹性本构形式\cite{flaschel2021hyperelastic}。本构人工神经网络通过输入不变量和网络组合表示应变能及材料响应\cite{linka2021constitutive,linka2023cann}。基于形态响应的参数反演通常先确定本构形式，再预测其中一个或多个参数，其输出表示相应本构模型和观测条件下的等效材料参数。

轮廓闭合、面积、周长和质心运动等关系可用于约束形态重构及参数预测。闭合轮廓需要保持首尾连续，面积和周长应与预测边界一致，质心位置与速度应符合时间顺序。几何关系和运动学关系能够提高不同输出之间的一致性，控制方程残差则属于更强的物理约束层级。二者在模型设计和结果解释中需要分别表述。

\subsection{小样本、噪声与多参数反演}

数值模拟可以生成连续时间帧和多个派生样本，独立样本数量仍由完整物理过程数量决定。同一计算过程产生的不同时间帧、稀疏组合和噪声版本共享相同材料参数和动力学轨迹。以完整过程划分训练集、验证集和测试集能够减少相关样本跨集合分布造成的信息泄漏\cite{roberts2017crossvalidation}。

小样本条件下，高维图像和多帧轮廓模型容易出现训练误差较低而独立测试误差较高的情况。单个轮廓可包含大量坐标，多个时间帧拼接后输入维度继续增加，但相邻轮廓点和相邻时间帧具有较强相关性。低维轮廓表示、代表性阶段采样和与数据规模相匹配的模型复杂度能够减少冗余信息\cite{vabalas2019limited}。观测噪声主要来自成像、分割、轮廓采样和时间定位，不同噪声形式会改变不同层级的形态信息。

多参数反演面临参数混淆。弹性模量和下降速度均可改变最大变形，表面张力和接触角均可改变铺展与回缩，细胞尺寸还会改变液体间隙和几何阻塞程度。内部形态、外部形态和不同动力学阶段的联合输入能够增加参数区分信息，其作用需要通过单参数任务、多参数任务和不同输入组合分别评价\cite{raue2009identifiability}。随机划分主要评价给定参数域内的插值能力，连续参数区间留出、参数组合留出和范围边缘测试则对应不同层级的泛化评价\cite{roberts2017crossvalidation}。

数值数据与实验数据之间还存在观测域差异。数值模型通常采用理想几何、确定物性和受控边界条件，实验图像则包含视角、景深、光照、细胞个体差异和设备噪声。数值数据适合检验形态信息、参数敏感性和反演方法，实验应用还需完成观测方式匹配、噪声建模和独立标定\cite{karniadakis2021piml,ovadia2019uncertainty}。

% ------------------------------------------------------------------
% 以下两节按当前工作安排冻结。待第三章至第五章完成后统一修订。
% ------------------------------------------------------------------
\section{现有研究不足与本文研究思路}

现有细胞力学测量、微流控变形分析、复合液滴动力学和数据驱动参数预测已形成较完整的方法体系。传统测量提供明确载荷和参数解释，微流控方法提高形态数据获取通量，数值模拟给出材料参数与动态响应之间的正向关系，机器学习用于表示复杂形态与参数之间的非线性映射。本文研究涉及的主要不足集中在观测完整性、形态表示、参数可辨识性和模型可靠性四个方面。

第一，稀疏形态观测中的力学信息保持缺少统一评价。现有微流控参数估计通常使用单帧几何指标、指定时刻图像或较完整的形态序列。实际成像可能只覆盖少量位置和时刻，轮廓分割还会引入局部边界误差。稀疏到完整的轮廓恢复首先是几何重构问题，同时也是力学信息恢复问题。平均点坐标误差较低的轮廓可能改变面积、周长、曲率或关键变形阶段，从而影响弹性模量预测。现有评价较少同时覆盖局部轮廓、整体几何、质心运动和下游参数识别。

本文将微通道细胞变形过程划分为多个空间和时间阶段，以少量观测轮廓恢复完整过程。点坐标表示用于评价局部边界误差，面积、圆度和长宽比用于评价整体形态，固定参数预测模型用于评价重构结果中的弹性模量信息。数据按照完整数值过程划分，避免同一轨迹的不同时间帧跨越训练集和测试集。含噪输入与无噪输入分别用于检验几何稳定性和力学信息保持。

第二，高维轮廓数据的低维表示和动态信息利用仍不充分。离散点坐标能够完整记录边界，但输入维度高，对轮廓起点、点序和局部噪声较敏感。面积、圆度和变形指数等低维指标具有明确含义，却无法描述局部边界和形态方向。完整时间序列还包含细胞位置和速度信息，单一轮廓表示难以同时描述形态变化和整体运动。

本文采用傅里叶轮廓表示压缩闭合边界，并由傅里叶系数计算形态派生量。质心位置和速度作为独立动态变量进入联合预测任务。不同傅里叶阶数用于比较低维表示能力和局部细节保持，模型结构、残差分支和损失项通过消融实验评价。该部分与点坐标重构形成递进关系：前者检验局部边界从稀疏观测到完整过程的恢复，后者检验低维形态、运动和弹性模量的联合预测。

第三，含细胞复合液滴中的多参数响应存在混淆。外部液滴铺展同时受下降速度、表面张力、黏度和接触角影响，内部细胞变形同时受流体载荷、材料参数、尺寸和空间位置影响。某一参数变化能够改变响应，不代表该响应可唯一确定该参数。弹性模量变化和下降速度变化可能产生相似的最大变形，表面张力变化和接触角变化可能产生相似的外部铺展。不同参数还可能在早期和晚期表现出不同敏感性。

本文从标准工况出发，分别比较材料、流动、界面和几何参数变化下的内部形态、质心运动、应力、应变能以及外部铺展和动态润湿响应。分析重点由单条曲线的变化方向扩展到响应幅度、峰值时刻、恢复过程和内外响应之间的关系。参数--响应敏感性用于筛选具有区分能力的观测量，响应相关性和相近形态用于识别参数混淆。第四章给出正问题和可辨识性基础，第五章利用相同参数体系检验由动态形态反演参数的稳定性。

第四，小样本条件下的多参数反演、泛化和不确定性证据不足。多参数数据库的规模受数值计算成本限制，高维图像或多帧轮廓模型容易在有限独立过程中产生过拟合。随机划分主要评价已覆盖参数域内的插值，不能替代未见参数组合和范围边缘测试。模型平均误差也无法反映参数混淆区和高误差样本。

本文以完整计算过程为独立样本，明确训练集、验证集和测试集的参数分布。单参数和多参数任务分别评价，图像、轮廓和几何指标作为不同输入形式进行比较。帧数组合、内部与外部形态、模型结构和损失项通过消融实验分析。噪声输入、参数组合留出和范围边缘测试用于评价稳健性，重复划分或多随机种子结果用于估计训练不确定性。预测结论限定于本文数值模型、参数范围和观测方式。

第五，数值形态与实际观测之间存在模型和观测域差异。微通道与复合液滴数值模型采用确定的材料关系、边界条件和几何设置，真实细胞具有结构异质性和个体差异。数值轮廓可直接由计算边界提取，实验轮廓还受到成像分辨率、投影方向、光学模糊和分割算法影响。数值数据库适合检验方法、响应敏感性和参数可辨识性，其结果不能直接替代实验标定。

本文在数值域内统一形态提取、坐标对齐、采样和归一化规则，并通过人工扰动检验观测误差影响。第三章重点处理稀疏帧和轮廓噪声，第四章明确各响应量的后处理定义，第五章比较不同观测输入对参数预测的影响。实验迁移、真实细胞结构和跨设备泛化作为当前工作的适用边界，不在数值结果中作超出证据的定量推断。

基于上述问题，本文按照“形态恢复--响应分析--参数反演”的顺序组织研究。微通道研究以受控几何和压力加载为条件，检验少量轮廓能否恢复完整形态，并由重构结果估计弹性模量。傅里叶表示将局部点坐标问题扩展为低维形态、质心运动和弹性模量的联合预测。含细胞液滴研究以多相界面和瞬态流固耦合为条件，分析多类参数对内外动态形态的影响，并检验不同参数的可辨识性。动态形态参数预测在此基础上比较单参数、多参数、不同输入和不同测试分布下的误差。

\section{本文主要研究内容与章节安排}

本文研究外部流体载荷作用下细胞动态形态与力学参数之间的关系。微通道场景采用压力驱动和几何受限加载，含细胞液滴场景采用下降接触产生的瞬态惯性、自由界面和润湿加载。两类场景均由数值模拟生成形态数据，研究内容分别面向稀疏形态恢复、低维动态表示、多参数响应分析和力学参数预测。

本文的第一项研究内容为微通道稀疏轮廓重构与弹性模量信息保持。变截面微通道数值模型生成细胞通过收缩区域的动态轮廓。完整过程按照空间位置和时间顺序进行对齐，少量观测时刻构成稀疏输入。MLP、VAE 和 Transformer 等模型用于恢复未观测时刻的轮廓。结果评价包括边界点误差、面积、周长、圆度、长宽比和噪声稳定性。重构轮廓随后输入固定弹性模量预测模型，用于判断几何接近是否同时对应材料参数信息保持。

第二项研究内容为基于傅里叶轮廓表示的形态、质心运动与弹性模量联合预测。细胞闭合边界转换为有限阶傅里叶系数，质心位置和速度作为独立运动变量。联合模型由稀疏观测预测完整形态过程、质心速度和弹性模量。傅里叶阶数分析用于确定低维表示与局部细节之间的平衡，模型比较和消融实验用于评价网络结构、残差分支和损失项的作用。不同通道位置的误差和几何派生量用于分析模型在进入、压缩和恢复阶段的表现。

第三项研究内容为含细胞复合液滴下降接触过程的动力学及参数敏感性。数值模型同时描述外部液滴界面、壁面润湿、内部流体和细胞变形。标准工况用于划分接触、铺展、最大变形和回缩恢复等阶段。弹性模量、黏弹性特征时间、表面张力、下降速度、内部尺寸和壁面接触角等参数分别改变，内部细胞形态、质心运动、应力、应变能以及外部液滴铺展和动态接触角同步比较。综合分析给出不同响应对参数的敏感阶段和参数混淆关系。

第四项研究内容为基于复合液滴动态形态的力学参数预测。第四章生成的内外形态序列构成参数数据库，细胞轮廓、液滴轮廓、几何指标和多帧图像作为候选输入。单参数预测用于评价受控条件下的映射精度，多参数联合预测用于检验参数同时变化时的可辨识性。不同输入形式、帧数、模型和损失设置进行对比，噪声、参数组合留出和范围边缘测试用于评价稳健性与泛化边界。预测结果按照物理单位报告，并结合误差分布和高误差样本分析参数混淆。

全文共分为六章，各章内容如下。

第一章为绪论。本章介绍细胞力学参数与形态响应的研究背景，梳理传统单细胞力学测量、微流控高通量形态测量、微通道参数反演、含细胞复合液滴下降接触、数值模拟以及人工智能参数预测的研究进展，归纳稀疏观测、形态表示、参数可辨识性、小样本和数值--实验域差异等问题，并给出本文的研究内容。

第二章为数学模型与研究方法。本章给出两类研究场景采用的流体控制方程、自由界面模型、细胞材料模型、流固耦合条件和 ALE 有限元方法，统一定义形态与运动指标、数据预处理方式、机器学习任务和评价指标。具体几何、参数范围、网格、时间步和网络设置分别在对应研究章节中给出。

第三章为基于稀疏形态响应的微通道细胞力学参数反演方法。本章建立变截面微通道细胞变形数据库，研究点坐标条件下的稀疏轮廓重构和噪声稳定性，并通过弹性模量预测评价重构轮廓的力学信息。傅里叶轮廓表示用于形态、质心运动和弹性模量的联合预测，模型结果通过阶数分析、基线比较、结构消融和区域误差进行评价。

第四章为含细胞复合液滴下降接触过程的数值模拟与力学响应分析。本章建立多相流动--可变形结构耦合模型，完成自由界面、流固耦合、网格和时间步相关验证。标准工况给出内外形态的阶段性演化，不同材料、流动、界面和几何参数的影响通过内部变形、质心运动、应力、应变能、铺展因子和动态接触角进行比较。综合敏感性和可辨识性分析用于确定第五章的输入响应和目标参数。

第五章为基于含细胞复合液滴动态形态的力学参数智能预测。本章定义单参数和多参数反演任务，建立多帧图像、轮廓和几何特征输入，给出模型结构、训练设置和数据划分。预测结果包括逐参数误差、输入和帧数消融、噪声稳定性、未见参数组合、参数混淆和不确定性分析。

第六章为全文总结与展望。本章按照稀疏形态恢复、动态响应分析和力学参数反演三个问题总结主要结果，归纳本文工作的创新内容，说明数值模型、样本规模、实验迁移和多参数可辨识性的限制，并讨论实验验证、三维模型、真实细胞结构、域适应和不确定性量化等后续研究方向。
